\section*{3.2 Author warm-up: A reduction from $r$-Bounded Set Cover to MDCP}
We will use instances of MDCP with a particular structure. This makes the second reduction from MDCP to Partial-MCSP more transparent.
Let $V$ be a ground set, and $\mathcal B=\{B_1,\ldots,B_n\}$ be a collection of nonempty subsets of $V$. For an element $v\in V$, we use $v^\uparrow\in\{0,1\}^n$ (the ``lifted'' version of $v$) to denote the string with the property that $v_i^\uparrow=1$ if and only if $v\in B_i$. We say that a collection $\mathcal B$ is complete (with respect to $V$) if the following holds: if $a,b\in V$ and $a^\uparrow=b^\uparrow$ then $a=b$. In other words, distinct elements of $V$ have different liftings with respect to $\mathcal B$. Our reduction from $r$-Bounded Set Cover to MDCP will always produce a family of generators that is complete.
Let $r$ be a large enough constant, so that say 10-approximating $r$-Bounded Set Cover is NP-hard. Given an instance $(1^n,\mathcal S)$ of this problem, the reduction proceed as follows. Fix $V\stackrel{\mathrm{def}}=[n^3]$. Partition $V$ into $n$ blocks $V_1,\ldots,V_n$, where $|V_i|=n^2$ for each $i\in[n]$. Given a set $A\subseteq V$, we let $A_i\stackrel{\mathrm{def}}=A\cap V_i$. Now view $V$ as the set $\{0,1\}^{3\log n}$, and for each $j\in[3\log n]$, let $B_j=\{v\in V\mid v_j=1\}$. For convenience, we let $\mathcal F=\{B_1,\ldots,B_{3\log n}\}$. Note that any family $\mathcal B$ of generators that contains $\mathcal F$ is complete with respect to $V$.
Let $A_{-i}\stackrel{\mathrm{def}}=\{A_1,\ldots,A_{i-1},A_{i+1},\ldots,A_n\}$. We say that a set $A$ is critical if, for every $i\in[n]$,
\[D(A\mid\mathcal F\cup A_{-i})>n\cdot\log n.\]
It is not hard to show that a uniformly random set $A\subseteq_{1/2}V$ is typically critical.
\setcounter{nttheorem}{8}
\begin{ntlemma}[Author Lemma9]
Let $A\subseteq_{1/2}V$ be sampled by letting $v\in A$ independently with probability$1/2$ for each $v\in V$. Then, $\Pr_A[A\text{ is critical}]\to1$ as $n\to\infty$.
\end{ntlemma}
\begin{proof}
For a fixed $i\in[n]$, we argue below that
\[\Pr_A[D(A\mid\mathcal F\cup A_{-i})\leq n\cdot\log n]=o(1/n).\]
The lemma follows by a union bound.
Note that, conditioning on the choice of $A_1,\ldots,A_{i-1},A_{i+1},\ldots,A_n$, the set $A_i$ is still a uniformly distributed subset of $V_i$. Moreover, after we fix $\mathcal F$ and $A_{-i}$, any construction of a set $E\subseteq V$ from $\mathcal F\cup A_{-i}$ using $s$ operations can be described by a binary string of length at most $O(s\cdot(\log s+\log n))$. Since $|V_i|=n^2$, the probability that the discrete complexity of $A$ (conditioned on the choice of $A_1,\ldots,A_{i-1},A_{i+1},\ldots,A_n$) given $\mathcal F\cup A_{-i}$ falls below $s=n\cdot\log n$ is at most
\[\frac{2^{O(s\cdot(\log s+\log n))}}{2^{n^2}}=o(1/n).\]
Since this holds for any choice of $A_1,\ldots,A_{i-1},A_{i+1},\ldots,A_n$, the desired probability upper bound holds.
\end{proof}
We assume from now on that we have efficiently produced a critical set $A$. (Our randomized reduction will always be correct on a critical set.) Note that so far we have only inspected the input $1^n$ from $(1^n,\mathcal S)$. Our approximate reduction outputs a triple $(1^{n^3},A,\mathcal B_{\mathcal S})$, where $A$ is generated as above, and whose family $\mathcal B_{\mathcal S}$ is defined as follows. For each set $S\in\mathcal S$, let $W_S\stackrel{\mathrm{def}}=\bigcup_{i\in S}A_i$. Now set
\[\mathcal B_{\mathcal S}\stackrel{\mathrm{def}}=\mathcal F\cup\{W_S\mid S\in\mathcal S\}.\]
Clearly, the triple $(1^{n^3},A,\mathcal B_{\mathcal S})$ can be efficiently computed from $(1^n,\mathcal S)$. We argue next that computing $D(A\mid\mathcal B_{\mathcal S})$ allows us to 2-approximate $\operatorname{cover}([n],\mathcal S)$.
\begin{ntlemma}[Author Lemma10]
If $\operatorname{cover}([n],\mathcal S)=\ell$ then $D(A\mid\mathcal B_{\mathcal S})\leq\ell$.
\end{ntlemma}
\begin{proof}
Let $S_1,\ldots,S_\ell$ be a cover of $[n]$ using sets from $\mathcal S$. Then, by construction of the sets $W_S$, we get that $A=W_{S_1}\cup\cdots\cup W_{S_\ell}$. In particular, it is possible to construct $A$ using at most $\ell$ operations (unions) starting from the sets in $\mathcal B_{\mathcal S}$.
\end{proof}
\begin{ntlemma}[Author Lemma11]
If $D(A\mid\mathcal B_{\mathcal S})=\ell$ then $\operatorname{cover}([n],\mathcal S)\leq2\ell$.
\end{ntlemma}
\begin{proof}
Let $t\stackrel{\mathrm{def}}=|\mathcal S|$, and suppose that $E_1,\ldots,E_\ell$ represent a construction of $A=E_\ell$ from sets in $\mathcal B_{\mathcal S}$. For convenience, we write $\mathcal B_{\mathcal S}=\mathcal F\cup\{W_{S_1},\ldots,W_{S_t}\}$. In order to prove Lemma11, we need the following simple claim.
\begin{ntclaim}[Author Claim12]
$E_1,\ldots,E_\ell$ depend on at most $2\ell$ sets from $\{W_{S_1},\ldots,W_{S_t}\}$.
\end{ntclaim}
This claim holds simply because each set $E_j$ depends on at most 2 other sets. One may replace $2\ell$ by $\ell+1$ by arguing more carefully, that any $\ell$-node directed acyclic graph with fan-in 2 and a single sink node has at most $\ell+1$ source nodes. But this tighter bound is unnecessary for our purpose.
Continuing with the proof of Lemma11, and relabelling some indexes if necessary, we assume for convenience of notation that the construction of $A$ from $\mathcal B_{\mathcal S}$ depends only on $\mathcal F\cup\{W_{S_1},\ldots,W_{S_{2\ell}}\}$. In other words,
\[D(A\mid\mathcal F\cup\{W_{S_1},\ldots,W_{S_{2\ell}}\})\leq\ell.\]
Since $W_{S_j}=\bigcup_{i\in S_j}A_i$ and $|S_j|\leq r$, we get that $D(W_{S_j}\mid\{A_i\}_{i\in S_j})\leq r$. In turn, by composing constructions, we obtain that
\[D(A\mid\mathcal F\cup\{A_i\}_{i\in S_1}\cup\cdots\cup\{A_i\}_{i\in S_{2\ell}})\leq\ell\cdot r.\]
We can assume that $r<\log n$ for a large enough $n$. In addition, notice that if $\ell>n$ then the statement of Lemma11 is trivial, since by assumption $\mathcal S$ always covers $[n]$ and as a consequence $\operatorname{cover}([n],\mathcal S)\leq n$. But from $r<\log n$ and $\ell\leq n$ it follows that the expression above can be upper bounded by $n\cdot\log n$. Given that $A$ is critical, it must be the case that the sets $S_1,\ldots,S_{2\ell}$ cover $[n]$. In other words, $\operatorname{cover}([n],\mathcal S)\leq2\ell$, which completes the proof.
\end{proof}
Consequently, it immediately follows from Lemmas9,10,and11 and from Theorem8 that MDCP is NP-hard to compute under polynomial-time randomized reductions.
